\chapter{Strings: hashing, KMP, Z}\label{ch:strings}
\index{strings}\index{rolling hash}\index{prefix function}\index{KMP}\index{Z-function}\index{trie}\index{border}\index{period}

String problems in contests reduce to a handful of linear-time
machines. After this chapter you can compare two substrings in
$O(1)$ --- and know exactly when a hash deserves trust --- match a
pattern in $O(n+m)$ two different ways, mine borders and periods
out of the prefix function, and walk a trie for anything that is
``a dictionary of prefixes'' in disguise.

\section{Terminology}\label{sec:strterm}
\index{substring}\index{subsequence}\index{prefix}\index{suffix}
A \emph{substring} is a contiguous block; a \emph{subsequence}
keeps order but may skip characters. Prefixes and suffixes anchor
the two ends. A \emph{border} of $s$ is a proper prefix that is
also a suffix: \texttt{abcababcab} has borders \texttt{ab} and
\texttt{abcab}. A \emph{period} is a distance $p$ with
$s_i = s_{i+p}$ for every $i < n-p$: \texttt{abcabca} has period
$3$ (and, vacuously, $7$). These four definitions --- border,
period, prefix function, Z-box --- are the whole vocabulary of
this chapter.

\section{Polynomial rolling hash}\label{sec:rhash}
\index{rolling hash}\index{hash collision}
Treat $s$ as a base-$B$ number: $H(s) = \sum_i s_i B^{\,n-1-i}
\bmod p$. Prefix hashes $h_i = H(s_0\dots s_{i-1})$ satisfy
$h_{i+1} = h_i B + s_i$, so any substring is recovered in $O(1)$:
$H(s_l\dots s_{r-1}) = h_r - h_l B^{\,r-l}$. Figure~\ref{fig:rhash}
shows the equivalent sliding-window view.
\begin{figure}[ht]\centering
\begin{tikzpicture}[font=\tiny]
\foreach \i/\c in {0/b,1/a,2/n,3/a,4/n,5/a,6/x}{\node[draw=slate!50,minimum size=0.32cm] at (\i*0.36,0) {\c};}
\draw[decorate,decoration={brace,amplitude=4pt}] (0.9*0.36,0.28) -- (4.1*0.36,0.28) node[midway,above=4pt] {$h(\texttt{ana})$};
\node[anchor=west] at (0,-0.7) {$h(\texttt{ana}) = \texttt{a}B^2 + \texttt{n}B + \texttt{a}$};
\node[anchor=west] at (0,-1.1) {slide right one: subtract $\texttt{a}B^2$, multiply by $B$, add \texttt{n}};
\node[anchor=west] at (0,-1.6) {$h(\texttt{nan}) = \bigl(h(\texttt{ana}) - \texttt{a}B^2\bigr)B + \texttt{n}$ \quad (all mod $p$)};
\end{tikzpicture}

\caption{Rolling one window right. Notice the dropped term carries
the full weight $B^{k-1}$ --- that is the term the update must
subtract first.}
\label{fig:rhash}
\end{figure}
\lstinputlisting{code/rhash.h}

A hash is a \emph{fingerprint}, not a proof: unequal strings can
share one. With $q$ queries the birthday bound kicks in --- among
$q$ hashes you should expect a collision once
$q \approx \sqrt{2p}$, so a single $10^9$ modulus is already
shaky at $q \sim 5{\cdot}10^4$, well inside contest range.
Doubling to two independent moduli moves the failure point to
$\sim\sqrt{2 p_1 p_2}$; the demo below makes the numbers concrete.

\begin{example}[Where collisions live]\label{ex:collide}
How much does modulus size matter, and does a second modulus fix it?
\par\smallskip
\textbf{Strategy.} Brute-force small alphabets under a toy modulus
to exhibit a \emph{real} colliding pair, then count collisions on
random strings under a mid-size modulus --- and under the
two-modulus pair used by \texttt{RH}.
\par\smallskip
\begin{solution}
\lstinputlisting{code/hash_demo.cpp}
Real output:
\begin{lstlisting}[numbers=none]
mod 97, base 26: "aaa" and "adt" both hash to 24
same pair, mod 1e9+7:   703 vs 800
same pair, mod 998244353: 703 vs 800
8192 8-letter strings: 5 extra hashes collide mod 9726647 (birthday), 0 under the RH pair
\end{lstlisting}
\end{solution}
\end{example}
\begin{exercise}[medium]\label{exr:birthday}
$q$ substrings are hashed mod $10^9{+}7$. Using
$1 - e^{-q^2/2p}$, at what $q$ does a collision become more likely
than not? Why does a second modulus change the exponent's
denominator?
\end{exercise}

\section{The prefix function (KMP)}\label{sec:kmp}
\index{prefix function}\index{failure link}
$\pi_i$ is the longest proper border of $s_0\dots s_i$. Compute
left to right: to extend $s_i$, start at the current border
candidate $j = \pi_{i-1}$ and keep falling along failure links
($j \leftarrow \pi_{j-1}$) until $s_i = s_j$; that character either
extends the border or leaves $\pi_i = 0$. The scan is $O(n)$
because $j$ only decreases inside the while-loop and can decrease
at most as many times as it ever increased.
\lstinputlisting{code/kmp.h}
The failure links make $\pi$ a deterministic automaton on the
pattern: state $=$ characters matched so far, transitions $=$
extend or fall. Figure~\ref{fig:kmpauto} draws the machine.
\begin{figure}[ht]\centering
% fig72-kmpauto.tex -- Figure 14.x: KMP automaton for pattern ABABAC; state = characters matched so far; dashed = failure links.
% fstyle.tex -- shared TikZ styles for the figures of this book (palette from lifeos.sty).
\tikzset{%
  cell/.style={draw=ink, minimum width=7mm, minimum height=6.5mm, inner sep=0pt, font=\footnotesize, fill=white},
  cellh/.style={cell, fill=accentsoft, draw=accent},
  cella/.style={cell, fill=amberbg, draw=amber},
  cellg/.style={cell, fill=goodbg, draw=good},
  cellx/.style={cell, fill=graybg, draw=slate, text=slate},
  vtx/.style={circle, draw=ink, minimum size=6.5mm, inner sep=0pt, font=\footnotesize, fill=white},
  vtxh/.style={vtx, draw=accent, fill=accentsoft},
  vtxa/.style={vtx, draw=amber, fill=amberbg},
  vtxg/.style={vtx, draw=good, fill=goodbg},
  vtxd/.style={vtx, draw=slate, fill=graybg, text=slate},
  lab/.style={font=\footnotesize, text=slate},
  labA/.style={font=\footnotesize, text=accent},
  labM/.style={font=\footnotesize, text=amber},
  idx/.style={font=\tiny, text=slate},
  arr/.style={-{stealth}, thick, ink},
  arrA/.style={-{stealth}, thick, accent},
  arrM/.style={-{stealth}, thick, amber},
  arrG/.style={-{stealth}, thick, good},
  edge/.style={thick, ink},
  edgeA/.style={thick, accent},
  edgeM/.style={thick, amber},
}%
% \cellrow{prefix}{x0}{y}{style}{v0,v1,...}: a row of cells named prefix0, prefix1, ... (pitch 0.75cm)
\providecommand{\cellrow}[5]{\foreach \v [count=\i from 0] in {#5} {\node[#4] (#1\i) at (#2+\i*0.75,#3) {\v};}}%
% \cellidx{prefix}{n}{start}: index labels under cells prefix0 .. prefix(n-1)
\providecommand{\cellidx}[3]{\foreach \i in {0,...,\numexpr#2-1\relax} {\node[idx, below=0pt of #1\i] {\the\numexpr\i+#3\relax};}}%
% \fitfig{<tikzpicture>}: shrink the picture to \linewidth only when it would overflow.
\ifdefined\fitbox\else\newsavebox{\fitbox}\fi%
\providecommand{\fitfig}[1]{\sbox{\fitbox}{#1}\ifdim\wd\fitbox>\linewidth\resizebox{\linewidth}{!}{\usebox{\fitbox}}\else\usebox{\fitbox}\fi}%

\begin{tikzpicture}[font=\footnotesize]
\foreach \i in {0,...,5} \node[vtx] (q\i) at (\i*2.1,0) {\i};
\node[vtxg, double, double distance=1.2pt] (q6) at (12.6,0) {6};
\foreach \a/\b/\c in {0/1/A, 1/2/B, 2/3/A, 3/4/B, 4/5/A, 5/6/C} \draw[arr, shorten >=1pt] (q\a) -- (q\b) node[midway, above, labA] {\c};
\foreach \a/\b/\l in {2/0/0, 3/1/1, 4/2/2, 5/3/3} \draw[arrM, dashed, bend left=48, shorten >=1pt] (q\a) to node[below, labM, font=\scriptsize] {fail $\to$ \l} (q\b);
\draw[arrM, dashed, bend left=22, shorten >=1pt] (q6.south) to node[below, labM, font=\scriptsize, fill=white, inner sep=1pt] {fail $\to$ 0} (q0.south);
\node[lab, anchor=west, align=left] at (0,-3.7) {at state 5 reading B: \texttt{C} was expected, so fall back $5\to3$ (\texttt{ABA}), where B is expected: advance to state $4$};
\end{tikzpicture}

\caption{KMP automaton for \texttt{ABABAC} ($\pi =
0,0,1,2,3,0$). Notice a failure is not a restart --- state 5 on
input B falls to 3 and immediately re-matches to 4.}
\label{fig:kmpauto}
\end{figure}
\begin{exercise}[easy]\label{exr:pihand}
Compute $\pi$ by hand for \texttt{ababac}, then for
\texttt{aaaaaa}. What does the second tell you about worst-case
border structure?
\end{exercise}

\section{The Z-function}\label{sec:z}
\index{Z-function}\index{Z-box}
$z_i$ is the longest common prefix of $s$ and $s_i\dots s_{n-1}$
($z_0$ unused). The algorithm keeps a \emph{Z-box} $[l,r]$: the
rightmost window already proved to match a prefix. Inside the box,
$z_i \ge \min(r-i+1,\; z_{i-l})$ for free; outside, or when the
copy might run past $r$, extend by direct comparison and grow the
box. Each extension moves $r$ right, so the total work is $O(n)$.
\lstinputlisting{code/zfunc.h}
\begin{figure}[ht]\centering
\begin{tikzpicture}[font=\tiny]
\foreach \i/\c in {0/a,1/a,2/b,3/a,4/a,5/b,6/a,7/a,8/b,9/a,10/a,11/b}{\node[draw=slate!50,minimum size=0.3cm] at (\i*0.34,0) {\c};}
\foreach \i/\v in {1/1,2/0,3/9,4/1,5/0,6/6,7/1,8/0,9/3,10/1,11/0}{\node[draw=amber!60,minimum size=0.3cm,below=0.1cm] at (\i*0.34,-0.35) {\v};}
\draw[decorate,decoration={brace,amplitude=4pt}] (3*0.34-0.15,0.25) -- (11*0.34+0.15,0.25) node[midway,above=5pt] {Z-box $[l,r]=[3,11]$ from $z[3]=9$};
\node[anchor=west] at (0,-1.0) {at $i{=}6$: inside the box, $z[6]\ge\min(11{-}6{+}1,\,z[6{-}3])=\min(6,9)=6$ --- free};
\node[anchor=west] at (0,-1.4) {$r$ only moves right; total extension work $\le n$};
\end{tikzpicture}

\caption{Z-box on \texttt{aabaabaabaab}. Notice $z_6 = 6$ costs no
comparisons --- it is copied from $z_3$ inside the box $[3,11]$.}
\label{fig:zbox}
\end{figure}

\section{Tries}\label{sec:trie}
\index{trie}
A trie stores a dictionary so that each node \emph{is} a prefix:
walking an edge appends a character; terminal markers record which
prefixes are real words. Insert and walk are $O(\text{length})$,
and a walk from position $i$ enumerates every dictionary word that
starts at $i$ --- exactly what word-splitting DP needs.
\lstinputlisting{code/trie.h}
\begin{figure}[ht]\centering
% fig75-trie.tex -- Figure 14.x: trie of dict = {a, ab, abab, ababc, c, cb}; double circle = terminal.
% fstyle.tex -- shared TikZ styles for the figures of this book (palette from lifeos.sty).
\tikzset{%
  cell/.style={draw=ink, minimum width=7mm, minimum height=6.5mm, inner sep=0pt, font=\footnotesize, fill=white},
  cellh/.style={cell, fill=accentsoft, draw=accent},
  cella/.style={cell, fill=amberbg, draw=amber},
  cellg/.style={cell, fill=goodbg, draw=good},
  cellx/.style={cell, fill=graybg, draw=slate, text=slate},
  vtx/.style={circle, draw=ink, minimum size=6.5mm, inner sep=0pt, font=\footnotesize, fill=white},
  vtxh/.style={vtx, draw=accent, fill=accentsoft},
  vtxa/.style={vtx, draw=amber, fill=amberbg},
  vtxg/.style={vtx, draw=good, fill=goodbg},
  vtxd/.style={vtx, draw=slate, fill=graybg, text=slate},
  lab/.style={font=\footnotesize, text=slate},
  labA/.style={font=\footnotesize, text=accent},
  labM/.style={font=\footnotesize, text=amber},
  idx/.style={font=\tiny, text=slate},
  arr/.style={-{stealth}, thick, ink},
  arrA/.style={-{stealth}, thick, accent},
  arrM/.style={-{stealth}, thick, amber},
  arrG/.style={-{stealth}, thick, good},
  edge/.style={thick, ink},
  edgeA/.style={thick, accent},
  edgeM/.style={thick, amber},
}%
% \cellrow{prefix}{x0}{y}{style}{v0,v1,...}: a row of cells named prefix0, prefix1, ... (pitch 0.75cm)
\providecommand{\cellrow}[5]{\foreach \v [count=\i from 0] in {#5} {\node[#4] (#1\i) at (#2+\i*0.75,#3) {\v};}}%
% \cellidx{prefix}{n}{start}: index labels under cells prefix0 .. prefix(n-1)
\providecommand{\cellidx}[3]{\foreach \i in {0,...,\numexpr#2-1\relax} {\node[idx, below=0pt of #1\i] {\the\numexpr\i+#3\relax};}}%
% \fitfig{<tikzpicture>}: shrink the picture to \linewidth only when it would overflow.
\ifdefined\fitbox\else\newsavebox{\fitbox}\fi%
\providecommand{\fitfig}[1]{\sbox{\fitbox}{#1}\ifdim\wd\fitbox>\linewidth\resizebox{\linewidth}{!}{\usebox{\fitbox}}\else\usebox{\fitbox}\fi}%

\begin{tikzpicture}[font=\footnotesize, level distance=1.1cm]
\tikzset{tn/.style={circle, draw=ink, minimum size=7mm, inner sep=0pt, fill=white},
         tt/.style={circle, draw=good, double, double distance=1.2pt, minimum size=7mm, inner sep=0pt, fill=goodbg}}
\node[tn] (root) at (0,0) {$\cdot$};
\node[tt] (a)  at (-2.4,-1.2) {a};
\node[tt] (ab) at (-2.4,-2.4) {b};
\node[tn] (aba) at (-2.4,-3.6) {a};
\node[tt] (abab) at (-2.4,-4.8) {b};
\node[tt] (ababc) at (-2.4,-6.0) {c};
\node[tt] (c) at (2.4,-1.2) {c};
\node[tt] (cb) at (2.4,-2.4) {b};
\draw[edge] (root)--(a) (a)--(ab) (ab)--(aba) (aba)--(abab) (abab)--(ababc) (root)--(c) (c)--(cb);
\node[lab, anchor=west, align=left] at (4.3,-1.0) {double circle = a word ends here (count $>0$)};
\node[lab, anchor=west, align=left] at (4.3,-1.8) {words: \texttt{a, ab, abab, ababc, c, cb}};
\node[labM, anchor=west, align=left, text width=6.5cm] at (4.3,-3.2) {text \texttt{ab}: walk $\cdot\to a\to b$ and count terminals on the way: \texttt{a} and \texttt{ab}, so $+2$ matches};
\end{tikzpicture}

\caption{Trie over $\{\texttt{a, ab, abab, ababc, c, cb}\}$.
Notice terminals can sit mid-path: \texttt{ab} ends inside the
edge to \texttt{abab}, which is what makes a walk report all
matches, not just the longest.}
\label{fig:trie}
\end{figure}
\begin{example}[CSES \emph{Word Combinations}]\label{ex:words}
Count the ways to write $s$ as a sequence of dictionary words,
$\bmod\ 10^9{+}7$.
\par\smallskip
\textbf{Strategy.} $\mathrm{dp}_i$ = ways to split $s_{i\dots n}$.
From each $i$, walk the trie on $s_i, s_{i+1}, \dots$; every
terminal at depth $d$ contributes $\mathrm{dp}_{i+d}$. Store
\emph{counts} at terminals (the dictionary may repeat words).
\par\smallskip
\begin{solution}
\lstinputlisting{code/word_combinations.cpp}
Real output (CSES sample):
\begin{lstlisting}[numbers=none]
2
\end{lstlisting}
\end{solution}
\end{example}
The same structure on bits --- the \emph{xor-trie} --- answers
``max $a_i \oplus x$'' queries and is named for the 1900 route.

\section{Applications}\label{sec:strapp}
\textbf{Matching.} KMP and Z solve the same problem on the same
concatenation $p\#s$: an occurrence ends wherever $\pi_i = m$
(KMP) or starts wherever $z_i = m$ (Z).
\begin{example}[CSES \emph{String Matching}, twice]\label{ex:match}
Count occurrences of $p$ in $s$.
\par\smallskip
\textbf{Strategy.} Build both machines on $p\#s$ and print both
counts --- for the official problem, print either.
\par\smallskip
\begin{solution}
\lstinputlisting{code/string_matching.cpp}
Real output (CSES sample: both methods, one line):
\begin{lstlisting}[numbers=none]
2 2
\end{lstlisting}
\end{solution}
\end{example}

\textbf{Borders.} $\pi_{n-1}$ is the longest border; iterating
$j \leftarrow \pi_{j-1}$ enumerates them all, longest first.
\begin{example}[CSES \emph{Finding Borders}]\label{ex:borders}
List all border lengths of $s$.
\par\smallskip
\textbf{Strategy.} Walk the $\pi$-chain from $\pi_{n-1}$, collect,
sort ascending.
\par\smallskip
\begin{solution}
\lstinputlisting{code/borders.cpp}
Real output (CSES sample):
\begin{lstlisting}[numbers=none]
2 5
\end{lstlisting}
\end{solution}
\end{example}
\begin{figure}[ht]\centering
% fig76-borders.tex -- Figure 14.x: borders of s = abcababcab (n = 10): longest border abcab (5); its border ab (2).
% fstyle.tex -- shared TikZ styles for the figures of this book (palette from lifeos.sty).
\tikzset{%
  cell/.style={draw=ink, minimum width=7mm, minimum height=6.5mm, inner sep=0pt, font=\footnotesize, fill=white},
  cellh/.style={cell, fill=accentsoft, draw=accent},
  cella/.style={cell, fill=amberbg, draw=amber},
  cellg/.style={cell, fill=goodbg, draw=good},
  cellx/.style={cell, fill=graybg, draw=slate, text=slate},
  vtx/.style={circle, draw=ink, minimum size=6.5mm, inner sep=0pt, font=\footnotesize, fill=white},
  vtxh/.style={vtx, draw=accent, fill=accentsoft},
  vtxa/.style={vtx, draw=amber, fill=amberbg},
  vtxg/.style={vtx, draw=good, fill=goodbg},
  vtxd/.style={vtx, draw=slate, fill=graybg, text=slate},
  lab/.style={font=\footnotesize, text=slate},
  labA/.style={font=\footnotesize, text=accent},
  labM/.style={font=\footnotesize, text=amber},
  idx/.style={font=\tiny, text=slate},
  arr/.style={-{stealth}, thick, ink},
  arrA/.style={-{stealth}, thick, accent},
  arrM/.style={-{stealth}, thick, amber},
  arrG/.style={-{stealth}, thick, good},
  edge/.style={thick, ink},
  edgeA/.style={thick, accent},
  edgeM/.style={thick, amber},
}%
% \cellrow{prefix}{x0}{y}{style}{v0,v1,...}: a row of cells named prefix0, prefix1, ... (pitch 0.75cm)
\providecommand{\cellrow}[5]{\foreach \v [count=\i from 0] in {#5} {\node[#4] (#1\i) at (#2+\i*0.75,#3) {\v};}}%
% \cellidx{prefix}{n}{start}: index labels under cells prefix0 .. prefix(n-1)
\providecommand{\cellidx}[3]{\foreach \i in {0,...,\numexpr#2-1\relax} {\node[idx, below=0pt of #1\i] {\the\numexpr\i+#3\relax};}}%
% \fitfig{<tikzpicture>}: shrink the picture to \linewidth only when it would overflow.
\ifdefined\fitbox\else\newsavebox{\fitbox}\fi%
\providecommand{\fitfig}[1]{\sbox{\fitbox}{#1}\ifdim\wd\fitbox>\linewidth\resizebox{\linewidth}{!}{\usebox{\fitbox}}\else\usebox{\fitbox}\fi}%

\begin{tikzpicture}[font=\footnotesize]
\foreach \c [count=\i from 0] in {a,b,c,a,b,a,b,c,a,b}
  \node[\ifnum\i<5 cellh\else cella\fi, minimum width=7mm] (s\i) at (\i*0.78,0) {\c};
\draw[decorate,decoration={brace,raise=3pt},accent,thick] (s0.north west) -- (s4.north east) node[midway,above=5pt,labA]{prefix \texttt{abcab}};
\draw[decorate,decoration={brace,mirror,raise=3pt},amber,thick] (s5.south west) -- (s9.south east) node[midway,below=5pt,labM]{suffix \texttt{abcab}: border of length $5$};
\draw[good, very thick] (s0.north west) rectangle (s1.south east); \draw[good, very thick] (s8.north west) rectangle (s9.south east);
\node[font=\footnotesize, text=good, anchor=north] at (s9.south |- 0,-1.55) {the inner border \texttt{ab} (length $2$) is outlined in green};
\node[lab, anchor=west] at (0,-2.4) {chain of borders: $\pi_{10}=5$, then the border of \texttt{abcab} is $\pi_5=2$, then $\pi_2=0$};
\end{tikzpicture}

\caption{Border chain on \texttt{abcababcab}. Notice each border
is itself a string whose own longest border is the next link ---
that is exactly the $\pi$-chain recursion.}
\label{fig:borders}
\end{figure}

\textbf{Periods.} $p$ is a period iff $n-p$ is a border
(Figure~\ref{fig:period}): the constraint $s_i = s_{i+p}$ for all
$i$ is precisely ``prefix of length $n-p$ equals suffix''. So the
same $\pi$-chain answers \emph{Finding Periods}, including
$p = n$ at the end.
\begin{figure}[ht]\centering
\begin{tikzpicture}[font=\tiny]
% s = abcabca, period 3: s[i]=s[i+3]; period = n - border
\foreach \i/\c in {0/a,1/b,2/c,3/a,4/b,5/c,6/a}{\node[draw=slate!50,minimum size=0.32cm] at (\i*0.36,0) {\c};}
\foreach \i/\j in {0/3,1/4,2/5,3/6}{\draw[<->,amber] (\i*0.36,0.32) -- (\j*0.36,0.32);}
\node[anchor=west] at (0,-0.6) {period $p{=}3$: $s_i = s_{i+p}$; suffix ``abca'' repeats prefix ``abca''};
\node[anchor=west] at (0,-1.0) {$\Leftarrow$ border of length $n-p = 4$; every border gives a period};
\end{tikzpicture}

\caption{Border $\Leftrightarrow$ period duality. Notice the
overlap: a period $p$ forces the prefix and suffix of length
$n-p$ to coincide.}
\label{fig:period}
\end{figure}
\begin{example}[CSES \emph{Finding Periods}]\label{ex:periods}
List all period lengths of $s$.
\par\smallskip
\textbf{Strategy.} Same $\pi$-chain as Finding Borders, printing
$n-j$ instead of $j$; append $n$.
\par\smallskip
\begin{solution}
\lstinputlisting{code/periods.cpp}
Real output (CSES sample):
\begin{lstlisting}[numbers=none]
3 6 7
\end{lstlisting}
\end{solution}
\end{example}
\begin{exercise}[medium]\label{exr:periodid}
Prove the identity used above: $p$ is a period of $s$ iff $n-p$
is a border. Conclude the smallest period is $n - \pi_{n-1}$.
\end{exercise}

\textbf{Substring comparison.} With prefix hashes, the LCP of any
two suffixes is one binary search away (compare $\ge$-mid
hashes), giving lexicographic comparison of any two substrings in
$O(\log n)$ --- no suffix array needed.
\begin{example}[CSES \emph{Minimal Rotation}]\label{ex:rotation}
Print the lexicographically smallest rotation of $s$.
\par\smallskip
\textbf{Strategy.} Rotations are length-$n$ windows of $s{+}s$.
Compare candidates with \texttt{RH::lcp} on $s{+}s$: jump past the
common prefix and decide on the first differing character ---
$O(\log n)$ per comparison.
\par\smallskip
\begin{solution}
\lstinputlisting{code/minimal_rotation.cpp}
Real output (CSES sample):
\begin{lstlisting}[numbers=none]
abac
\end{lstlisting}
\end{solution}
\end{example}

\begin{table}[ht]\centering\small
\begin{tabularx}{\linewidth}{l l l X}
\toprule
task & tool & cost & gotcha\\
\midrule
substring equality & prefix hashes, two moduli & $O(1)$/q after $O(n)$ & birthday bound at $\sqrt{2p}$\\
pattern matching & KMP or Z on $p\#s$ & $O(n+m)$ & sentinel must not appear in alphabet\\
all borders & $\pi$-chain from $\pi_{n-1}$ & $O(n)$ & iterate $j\leftarrow\pi_{j-1}$, not $\pi_j$\\
all periods & same chain, print $n-j$ & $O(n)$ & append $p=n$ manually\\
dict word splits & trie walk per position & $O(n\cdot L)$ & terminals mid-path count too\\
substring compare & hash $+$ binary-search LCP & $O(\log n)$/cmp & half-open $[l,r)$ in \texttt{get}\\
\bottomrule
\end{tabularx}
\caption{The task names the machine; the machine names the cost.}
\label{tab:strings}
\end{table}

\section{Common mistakes}
\index{common mistakes}
\begin{itemize}
\item \textbf{Skipping the sentinel.} Matching on $p+s$ (no
\texttt{\#}) lets a ``match'' straddle the boundary --- a
separator outside the alphabet is required.
\item \textbf{Wrong index in the $\pi$-chain.} The next candidate
border is $\pi_{j-1}$, not $\pi_j$; the off-by-one yields either
an infinite loop or the trivial answer.
\item \textbf{Negative residues in \texttt{get}.} $h_r - h_l
B^{r-l} \bmod p$ needs the $+p$ before returning --- C++ \texttt{\%}
keeps the sign of the dividend.
\item \textbf{One modulus, many queries.} Birthday collisions
arrive at $\sqrt{2p}$; on $q \gtrsim 10^5$ a single $10^9$ prime
is a gamble, not a proof.
\item \textbf{Assuming terminals only at leaves.} In word-splitting
tries, \texttt{cnt} must be read at \emph{every} node on the walk,
or the short matches vanish.
\end{itemize}

\section{Additional exercises}
\begin{exercise}[easy] CSES \emph{Finding Borders}. \end{exercise}
\begin{exercise}[easy] CSES \emph{Finding Periods}. \end{exercise}
\begin{exercise}[easy] USACO-B \emph{Censoring}. \end{exercise}
\begin{exercise}[easy] CF \emph{Check Transcription}. \end{exercise}
\begin{exercise}[easy] CSES \emph{Counting Patterns} --- $k$ independent KMP runs. \end{exercise}
\begin{exercise}[medium] CSES \emph{String Matching}. \end{exercise}
\begin{exercise}[medium] CSES \emph{Minimal Rotation}. \end{exercise}
\begin{exercise}[medium] CSES \emph{Longest Palindrome}. \end{exercise}
\begin{exercise}[medium] CSES \emph{Repeating Substring}. \end{exercise}
\begin{exercise}[medium] USACO-G \emph{Lights Out}. \end{exercise}
\begin{exercise}[medium] USACO-S \emph{Censoring (Silver)} --- stack + KMP. \end{exercise}
\begin{exercise}[hard] CSES \emph{Required Substring}. \end{exercise}
\begin{exercise}[hard] CSES \emph{Distinct Substrings} --- suffix structures or hashing + sort. \end{exercise}
\begin{exercise}[hard] CSES \emph{String Functions} --- compute both $\pi$ and $z$ sums. \end{exercise}
\begin{exercise}[hard] CSES \emph{Palindrome Queries} --- hashing both directions. \end{exercise}

\chapternotes

\section*{Further reading}
\begin{itemize}
\item A.~Laaksonen, \emph{Competitive Programmer's Handbook},
\url{https://cses.fi/book/book.pdf} --- free PDF; the
trie/hashing/Z treatment this chapter extends.
\item J.~Sannemo, \emph{Principles of Algorithmic Problem Solving}
(2025), \url{https://usaco.guide/PAPS1.pdf} --- free PDF; the
string chapter behind our proof style.
\item cp-algorithms, \emph{String hashing} --- collision bounds
and double-hash practice.
\newline\url{https://cp-algorithms.com/string/string-hashing.html}
\item cp-algorithms, \emph{Prefix function} --- the automaton
view and its applications catalogue.
\newline\url{https://cp-algorithms.com/string/prefix-function.html}
\item cp-algorithms, \emph{Z-function} --- the Z-box invariant
proved cleanly.
\newline\url{https://cp-algorithms.com/string/z-function.html}
\item The USACO Guide, \emph{Hashing} module,
\url{https://usaco.guide/gold/hashing} --- the contest settings
where hashing is the intended weapon.
\item D.~Gusfield, \emph{Algorithms on Strings, Trees, and
Sequences} (1997), ISBN 978-0521585194 --- the reference for
everything past this chapter: suffix trees, automata, and why
Z-values work.
\item S.~Halim, F.~Halim, S.~Effendy, \emph{Competitive
Programming 4} (2020), \url{https://cpbook.net/} --- the string
problem taxonomy.
\end{itemize}

\begin{rulebox}[Watch instead]
\begin{itemize}
\item \textbf{Abdul Bari --- ``9.1 Knuth-Morris-Pratt KMP String
Matching Algorithm''}, 18m55s,
\url{https://www.youtube.com/watch?v=V5-7GzOfADQ} --- covers
\S15.1, \S15.3: borders (he calls $\pi$ the ``LPS table'') and
the failure-link walk. (Verified 2026-10-11.)
\item \textbf{Abdul Bari --- ``9.2 Rabin-Karp String Matching
Algorithm''}, 23m50s,
\url{https://www.youtube.com/watch?v=qQ8vS2btsxI} --- covers
\S15.2: rolling hash as a sliding base-$B$ number; he uses
base-10 codes for pedagogy --- same idea, our modulus is prime.
(Verified 2026-10-11.)
\item \textbf{UC Davis --- ``Linear-time pattern matching:
Z-values and Z-algorithm''} (Dan Gusfield), 51m46s,
\url{https://www.youtube.com/watch?v=MFK0WYeVEag} --- covers
\S15.4: the Z-box invariant from the algorithm's original
expositor. (Verified 2026-10-11.)
\end{itemize}
\emph{Notation note:} Bari writes \texttt{lps} for our $\pi$ and
indexes strings from 1 in slides; Gusfield's $Z_k$, $l$, $r$
match the book exactly.
\end{rulebox}
Rolling-hash construction and collision practice follow
cp-algorithms and the USACO Guide~\cite{cpalgorithms,usacoguide};
the prefix-function, Z-function and trie treatments track
Laaksonen~\cite{laaksonen2024} and Sannemo~\cite{sannemo2025};
the Z-box proof shape follows Gusfield (see Further reading);
problem band and selection follow Halim~\cite{halimcp4} and the
official CSES set~\cite{cses}.
